\chapter{Geometric Ethics for engineers}\label{ch:ge}
% Author: agent A. Every claim carries a source comment. Paths are local clones.
% GE book = C:\source\geometric-ethics-content\source\Geometric-Ethics-v1.24.tex (chapter 8 at 10740-12230)
% GE updates = C:\source\geometric-ethics-content\source\updates\ (sec8_10b_executable_strata.tex, sec8_4_hohfeld_v4.tex)
% GE paper = C:\source\erisml-lib\docs\papers\foundations\Geometric_Ethics_Foundational_Paper.tex
% hohfeld.py (V4) = erisml-lib origin/main at commit df00fc5 (3 Oct 2026), read with git show; the
%   local working copy of erisml-lib still holds the earlier D4 file.
% MV reference = C:\source\erisml-lib\docs\moralvector_reference.md
% moral_vector.py = C:\source\erisml-lib\src\erisml\ethics\moral_vector.py

Geometric Ethics is the mathematical account of moral evaluation on which the ErisML stack is
built. This chapter gives an engineer the parts of it that the rest of the monograph uses. A moral
assessment has several independent dimensions, and reducing it to one number loses information
that cannot be recovered. The space of assessments is divided into strata, and the rules change
when a situation crosses from one stratum into another. The crossings happen at gates, and a gate
fires only on a declared kind of event. The four Hohfeldian positions of obligation, claim,
liberty and no-claim are related by a group of four operations. The software carries the
assessment as a vector of nine named dimensions. Each idea is stated with its status in the
sources, since some are proved, some are modeling axioms and some rest on measurement.

\section{A moral assessment has nine axes}\label{sec:ge-space}

Geometric Ethics models the morally relevant state of a system of agents as a point $p$ in a
space $\mathcal{M}$, the moral manifold. In the foundational paper $\mathcal{M}$ has nine
coordinate axes.
% src: GE paper:183-198 (Definition, Moral Manifold)
They are physical consequences, rights and duties, justice and fairness, autonomy, privacy,
societal impact, virtue and character, legitimacy, and epistemic status. The number nine is a
modeling choice whose adequacy is empirical. If a moral phenomenon appears that the nine axes
cannot represent, the space must be extended. If some axes prove redundant, it must be reduced.
% src: GE paper:201-203 (Remark, Modeling Choice)

The space carries a Riemannian metric $g_{\mu\nu}$, a positive-definite inner product on each
tangent space that prices a small displacement. The metric is not fixed. Its variation across
$\mathcal{M}$ encodes how the relative weight of the moral dimensions depends on context.
% src: C:\source\erisml-lib\docs\papers\foundations\Philosophy_Engineering_Foundation_v1.0.tex:313-322
For a path $\gamma\colon[0,1]\to\mathcal{M}$ through the space, the behavioural friction is the
path's length under that metric,
\begin{equation}
\mathrm{BF}[\gamma] = \int_0^1 \sqrt{g_{\mu\nu}(\gamma(t))\,\dot\gamma^{\mu}(t)\,\dot\gamma^{\nu}(t)}\;dt ,
\label{eq:ge-friction}
\end{equation}
with summation over repeated indices.
% src: Philosophy_Engineering_Foundation_v1.0.tex:324-331
The monograph needs only two consequences of this picture. Distances between assessments are
defined by a metric that can change with the situation. Paths through the space can be priced,
which is what the discrete decision complex of Section~\ref{sec:ge-discrete} does.

\section{A single number cannot carry the assessment}\label{sec:ge-scalar}

Let $T$ be a moral tensor of rank at least one, meaning an assessment with at least one index over
the moral dimensions. Let $c$ be any map that contracts $T$ to a scalar $S = c(T)$. Write
$I(A;B)$ for the mutual information between random variables $A$ and $B$, and treat the situation
being assessed as a random variable. The foundational paper's Information Monotonicity theorem
states
\begin{equation}
I(T;\,\text{situation}) \;\geq\; I(S;\,\text{situation}),
\label{eq:ge-infomono}
\end{equation}
with equality if and only if $T$ is already effectively scalar, so that all its dimensions are
collinear.
% src: GE paper:320-336 (Theorem, Information Monotonicity, and proof sketch)
The proof is the data processing inequality. A deterministic function of a random variable
cannot carry more information about a third variable than the variable itself carries.

The brochure states a second form, about structure instead of information.
% src: C:\Users\abptl\Documents\Philosophy-Engineering-Brochure.pdf p1 ("no continuous injection [0,1]^d->R for d>=2")
The Geometric Ethics manuscript proves it for nine dimensions as its Scalar Irrecoverability
theorem.
% src: C:\source\erisml-lib\docs\papers\foundations\ch20_economics.py:34 (Theorem 20.2 proof, invariance of dimension)

\medskip\noindent\textbf{Proposition (scalar collapse).} For $d \geq 2$ there is no continuous
injective map $\varphi\colon [0,1]^d \to \mathbb{R}$.

\medskip\noindent\textit{Proof.} The manuscript argues from Brouwer's invariance of dimension. A
shorter argument uses connectedness. Suppose $\varphi$ were continuous and injective. Its image is
compact and connected, hence a closed interval $[a,b]$, and $a<b$ because $\varphi$ is injective.
Pick $c$ with $a<c<b$ and let $q = \varphi^{-1}(c)$. For $d \geq 2$ the punctured cube
$[0,1]^d \setminus \{q\}$ is connected. Its image under $\varphi$ is $[a,b]\setminus\{c\}$, which
is not connected. A continuous map cannot send a connected set onto a disconnected one, so no
such $\varphi$ exists.\hfill$\square$
% src: proof by the chapter author (standard topology); the manuscript's proof is at ch20_economics.py:34

Any score that is continuous in the nine dimensions therefore sends distinct assessments to the
same number. Figure~\ref{fig:ge-scalar} shows the contraction and the brochure's example of a
collapsed verdict.

\begin{figure}[t]
\centering
% Scalar collapse: nine dimensions contracted to one number. No values are plotted.
% src: C:\source\erisml-lib\docs\moralvector_reference.md:16-30 (nine dimensions)
% src: C:\Users\abptl\Documents\Philosophy-Engineering-Brochure.pdf p1 ("0.83 safe" example)
% src: C:\source\erisml-lib\src\erisml\ethics\moral_vector.py:161-203 (to_scalar, weighted average)
\begin{tikzpicture}[house,
  dim/.style={detail, text width=19mm, minimum height=8mm}]
\node[panel=Blue, minimum width=68mm, minimum height=42mm] (v) at (0,0) {};
\node[stepname=Blue] at (0,1.75) {Moral vector, nine dimensions};
\node[dim] at (-2.25,0.9)  {physical harm};
\node[dim] at (0,0.9)      {rights};
\node[dim] at (2.25,0.9)   {fairness};
\node[dim] at (-2.25,-0.05) {autonomy};
\node[dim] at (0,-0.05)     {privacy};
\node[dim] at (2.25,-0.05)  {societal, environmental};
\node[dim] at (-2.25,-1.0) {virtue and care};
\node[dim] at (0,-1.0)     {legitimacy and trust};
\node[dim] at (2.25,-1.0)  {epistemic quality};
\node[panel=Orange, minimum width=30mm, minimum height=42mm] (s) at (5.7,0) {};
\node[stepname=Orange] at (5.7,1.75) {One number};
\node[desc] at (5.7,1.05) {$S=\sum_i w_i v_i$};
\node[font=\sffamily\bfseries\large, text=hsOrange] at (5.7,0.25) {0.83 safe};
\node[desc, text width=26mm] at (5.7,-0.85) {the brochure's example of a collapsed verdict};
\draw[flow] (v.east) -- (s.west);
\node[detail, draw=hsGreyLine, text width=70mm] at (2.85,-2.55) {the map from vector to number has no continuous inverse when $d\ge 2$};
\end{tikzpicture}

\caption{A weighted contraction of the nine dimensions to one number. The arrow runs one way
only, since distinct assessments land on the same number.}
\label{fig:ge-scalar}
\end{figure}

Neither result forbids a scalar at the moment of decision. The foundational paper defines the
satisfaction scalar as the contraction $S = I_\mu O^\mu$ of an interest covector $I_\mu$ with an
obligation vector $O^\mu$. It calls this contraction the moment when the tensor becomes a decision.
% src: GE paper:304-318 (Definition, Interest Covector; equation eq:satisfaction)
The engineering rule that follows is the brochure's. The scalar is the last step and not the
substrate.
% src: Brochure p1 ("Scalar is last step, not substrate")
DEME keeps both views. Its \code{MoralVector.to\_scalar} returns a weighted average for backward
compatibility, with default weights of 1.0 on six dimensions, 0.8 on societal and environmental
impact, 0.7 on virtue and care, and 0.5 on epistemic quality.
% src: moral_vector.py:28-41 (DEFAULT_DIMENSION_WEIGHTS), 161-203 (to_scalar)
Its \code{dominates} method compares two vectors without any weights. One vector dominates another
when it is at least as good on every dimension and strictly better on one.
% src: moral_vector.py:209-240

\section{Moral space is stratified}\label{sec:ge-strata}

A smooth manifold cannot hold a moral bright line. The foundational paper states that no smooth
manifold can represent at once continuous variation of moral weight inside one kind of situation,
a discontinuous change between kinds, and a drop in dimension at the boundary.
% src: GE paper:217-220 (Proposition, Insufficiency of Smooth Manifolds)
Geometric Ethics therefore uses a stratified space.

A stratification of a topological space $X$ is a locally finite partition $\{S_\alpha\}_{\alpha\in
A}$ of $X$ into smooth manifolds, called strata, of various dimensions, with a partial order
$\preccurlyeq$ on the index set $A$. It satisfies the frontier condition
\begin{equation}
S_\alpha \cap \overline{S_\beta} \neq \varnothing \;\Longrightarrow\; S_\alpha \subset
\overline{S_\beta} \ \text{and}\ \alpha \preccurlyeq \beta ,
\label{eq:ge-frontier}
\end{equation}
where $\overline{S_\beta}$ is the closure of $S_\beta$. Each point of $X$ has a neighbourhood that
meets only finitely many strata.
% src: GE book:10750-10766 (Definition 8.1)
Whitney's conditions A and B add regularity at the boundaries, so that tangent planes and secant
lines behave as a stratum approaches a lower one.\cite{whitney1965tangents}
% src: GE book:10774-10805 (Definitions 8.2, 8.3)
In moral terms the conditions make a boundary sharp and approachable. The evaluation changes
smoothly up to the boundary and jumps there.
% src: GE book:10811-10818

Section~8.3 of the manuscript sorts the boundaries into four types, shown in
Figure~\ref{fig:ge-boundaries}. Write $S$ for the satisfaction function on $\mathcal{M}$.
% src: GE book:10836-10847

\paragraph{Type I, thresholds.} A threshold boundary is a hypersurface $B = \phi^{-1}(0)$, the zero
set of a smooth function $\phi\colon\mathcal{M}\to\mathbb{R}$. On each side $S$ and $g_{\mu\nu}$ are
smooth. Across $B$ either may jump. The sign of $\phi$ says which side applies. The age of majority
and the threshold of informed consent are examples.
% src: GE book:10849-10900 (Definition 8.4 and examples table)

\paragraph{Type II, phase transitions.} At a phase transition boundary the metric $g_{\mu\nu}$
changes discontinuously in rank or signature. The set of trade-offs changes, not only the value
of $S$. The manuscript's example is a medical emergency, in which urgency takes lexicographic
priority.
% src: GE book:10909-10939 (Definition 8.5, emergency example)

\paragraph{Type III, absorbing strata.} A stratum $S_0$ is absorbing if it has a tubular
neighbourhood $N$ and a continuous retraction $r\colon N\to S_0$ such that every compatible
evaluation function factors through $r$,
\begin{equation}
S(p) = S(r(p)) \qquad \text{for all } p\in N .
\label{eq:ge-absorbing}
\end{equation}
Near an absorbing stratum the evaluation ignores the transverse coordinates. Absorbing strata
formalize nullifiers. The manuscript reports three that hold across all domains of the Dear Abby
corpus, namely abuse ($n=582$), danger ($n=218$) and impossibility ($n=144$).
% src: GE book:10970-11022 (Definition 8.6, nullifier table)

\paragraph{Type IV, constraint surfaces.} The constraint set $C\subset\mathcal{M}$ is closed, and
$S|_C = -\infty$. Its boundary $\partial C$ is the constraint surface. No finite benefit outweighs
crossing it.
% src: GE book:11039-11062 (Definition 8.7)

\begin{figure}[t]
\centering
% The four boundary types of Geometric Ethics section 8.3 and how the runtime realizes each.
% src: C:\source\geometric-ethics-content\source\Geometric-Ethics-v1.24.tex:10849-11070 (Defs 8.4-8.7)
% src: same file :12042-12055 (penalties beta_k by type, remark after Def 8.13)
% src: C:\source\geometric-ethics-content\source\updates\sec8_10b_executable_strata.tex:18-41
% src: C:\source\erisml-compiler-monitor\src\erisml_compiler\runtime\strata.py:43 (BOUNDARIES)
\begin{tikzpicture}[house,
  card/.style={panel=#1, minimum width=28mm, minimum height=60mm},
  dd/.style={desc, text width=25mm},
  rt/.style={detail, text width=24mm, minimum height=14mm}]
\node[card=Blue]   at (0,0)    {};
\node[card=Purple] at (3.15,0) {};
\node[card=Red]    at (6.3,0)  {};
\node[card=Grey]   at (9.45,0) {};
\node[step=Blue]   at (0,2.5)    {I};
\node[step=Purple] at (3.15,2.5) {II};
\node[step=Red]    at (6.3,2.5)  {III};
\node[step=Grey]   at (9.45,2.5) {IV};
\node[stepname=Blue]   at (0,1.85)    {Threshold};
\node[stepname=Purple] at (3.15,1.85) {Phase transition};
\node[stepname=Red]    at (6.3,1.85)  {Absorbing stratum};
\node[stepname=Grey]   at (9.45,1.85) {Constraint surface};
\node[dd] at (0,0.95)    {a level set $\phi^{-1}(0)$ where $S$ or $g$ may jump};
\node[dd] at (3.15,0.95) {the metric changes in rank or signature};
\node[dd] at (6.3,0.95)  {near it $S$ depends only on the retraction to the stratum};
\node[dd] at (9.45,0.95) {$S=-\infty$ on a closed forbidden set};
\node[dd] at (0,0.0)    {penalty $\beta$ finite};
\node[dd] at (3.15,0.0) {penalty $\beta$ large and finite};
\node[dd] at (6.3,0.0)  {penalty $\beta=\infty$};
\node[dd] at (9.45,0.0) {penalty $\beta=\infty$, no gate leads there};
\node[rt] at (0,-1.6)    {runtime gate marked \texttt{threshold}};
\node[rt] at (3.15,-1.6) {runtime gate marked \texttt{phase}};
\node[rt] at (6.3,-1.6)  {left only on a human-oversight event};
\node[rt] at (9.45,-1.6) {non-defeasible prohibition in the scene};
\end{tikzpicture}

\caption{The four boundary types of Geometric Ethics with the crossing penalty each carries and
the construct that realizes it at run time. Only the first two are gate kinds in the runtime,
since the third is a property of a stratum and the fourth is a prohibition.}
\label{fig:ge-boundaries}
\end{figure}

\section{Gates move a situation between strata}\label{sec:ge-gates}

Boundaries say where the rules change. Semantic gates say what makes them change.
% src: GE book:11077-11079

\medskip\noindent\textbf{Definition 8.8 (semantic gate).} A semantic gate is a map
\begin{equation}
G\colon \mathcal{F}\times\mathcal{S}_\alpha \to \mathcal{S}_\beta ,
\label{eq:ge-gate}
\end{equation}
where $\mathcal{F}$ is a space of triggering features, $\mathcal{S}_\alpha$ the source stratum and
$\mathcal{S}_\beta$ the target. The gate has a characteristic function
$\sigma\colon\mathcal{F}\to\{0,1\}$ and is defined only on $\sigma^{-1}(1)\times\mathcal{S}_\alpha$.
It fires when a feature with $\sigma(f)=1$ is present, and the evaluation then follows the rules
of $\mathcal{S}_\beta$.
% src: GE book:11081-11097

The defining property of a gate is that it is discrete. A transition is all or nothing, like a
step function and unlike a sigmoid. The manuscript treats this as an empirical claim. It reports
gate detection rates from the Dear Abby corpus between 89 and 97 percent, for triggers such as
``you promised'' and ``only if convenient'', with every listed gate behaving as a step function.
% src: GE book:11099-11145 (discreteness, gate detection table: 94, 89, 97, 92, 91 percent)

\medskip\noindent\textbf{Definition 8.9 (penumbral zone).} A penumbral zone around a boundary $B$
is the tubular neighbourhood
\begin{equation}
N_\varepsilon(B) = \{\, p\in\mathcal{M} \;\mid\; d(p,B) < \varepsilon \,\}
\label{eq:ge-penumbra}
\end{equation}
within which the governing stratum is uncertain. Here $d$ is the distance on $\mathcal{M}$ and
$\varepsilon>0$ is the zone's width.
% src: GE book:11365-11374
The manuscript stresses that the boundary itself stays sharp. What is open inside the zone is
which stratum governs a given point. Its example is capacity to consent, a boundary that is
definite but hard to measure in a particular person.
% src: GE book:11376-11396

\medskip\noindent\textbf{Definition 8.11 (boundary crossing data).} At a boundary $B$ between strata
$S_\alpha$ and $S_\beta$ the crossing data has three parts. The first is a transition map
$T_{\alpha\beta}$ from tensors on $S_\alpha$ at $B$ to tensors on $S_\beta$ at $B$. The second is
a jump function $\Delta S\colon B\to\mathbb{R}\cup\{-\infty\}$ that gives the discontinuity in
$S$. The third is a gate condition $G\colon B\to\{0,1\}$ that says whether the transition fires at
each point of $B$.
% src: GE book:11677-11690
The transition map need not be the identity. An obligation to keep a secret becomes an obligation
to disclose it when the secret turns out to involve harm to a third party. The map reverses the
obligation along the disclosure dimension.
% src: GE book:11692-11702

\section{Two discrete realizations, one for search and one for an agent}\label{sec:ge-discrete}

The stratified space is continuous. The manuscript gives it two discrete realizations, shown in
Figure~\ref{fig:ge-discrete}.
% src: GE book:11985-11992; GE updates sec8_10b_executable_strata.tex:9-14

\begin{figure}[t]
\centering
% From the continuous stratified space to the two discrete realizations of chapter 8.
% src: C:\source\geometric-ethics-content\source\Geometric-Ethics-v1.24.tex:11985-12075 (Defs 8.12, 8.13, Prop 8.5)
% src: C:\source\geometric-ethics-content\source\updates\sec8_10b_executable_strata.tex:7-37 (Def 8.14)
\begin{tikzpicture}[house,
  box/.style={panel=#1, minimum width=34mm, minimum height=36mm},
  dd/.style={desc, text width=30mm},
  ex/.style={detail, text width=29mm}]
\node[box=Purple] (m) at (0,0)   {};
\node[box=Blue]   (d) at (4.0,0) {};
\node[box=Teal]   (x) at (8.0,0) {};
\node[step=Purple] at (0,1.35)   {1};
\node[step=Blue]   at (4.0,1.35) {2};
\node[step=Teal]   at (8.0,1.35) {3};
\node[stepname=Purple] at (0,0.75)   {Stratified space};
\node[stepname=Blue]   at (4.0,0.75) {Decision complex};
\node[stepname=Teal]   at (8.0,0.75) {Executable strata};
\node[dd] at (0,0.1)   {smooth strata, metric $g_{\mu\nu}$, four boundary types};
\node[dd] at (4.0,0.1) {states, actions as weighted edges, for search};
\node[dd] at (8.0,0.1) {current stratum, gates fired by events, for an agent};
\node[ex] at (0,-1.05)   {Definitions 8.1 to 8.11};
\node[ex] at (4.0,-1.05) {$w=\Delta a^{\top}\Sigma^{-1}\Delta a+\sum_k\beta_k\,\mathbf{1}_k$};
\node[ex] at (8.0,-1.05) {$(\Sigma,\sigma_0,G,A,B)$};
\draw[flow] (m) -- (d);
\draw[flow] (m.north) to[bend left=22] node[above, flowlabel] {for an acting agent} (x.north);
\end{tikzpicture}

\caption{The continuous stratified space and its two discrete realizations. The decision complex
prices paths for search, and the executable stratification tracks one current stratum for an
agent that acts.}
\label{fig:ge-discrete}
\end{figure}

\paragraph{The decision complex.} Definition~8.12, a modeling axiom, builds a weighted simplicial
complex $\Delta$. Each vertex $v_i$ is a morally relevant state with a nine-dimensional attribute
vector $a(v_i)\in\mathbb{R}^9$. Each edge is an available action. Higher simplices are bundles of
actions taken together.
% src: GE book:11994-12006 (Definition 8.12)
Definition~8.13 prices an edge as
\begin{equation}
w(v_i,v_j) = \Delta a^{\top}\,\Sigma^{-1}\,\Delta a \;+\; \sum_k \beta_k\,
\mathbf{1}[\text{boundary } k \text{ crossed}] ,
\label{eq:ge-edge}
\end{equation}
where $\Delta a = a(v_j)-a(v_i)$, $\Sigma$ is the $9\times9$ moral covariance matrix, $\beta_k$ is
the penalty for crossing boundary $k$, and $\mathbf{1}[\cdot]$ is one when its condition holds and
zero otherwise.
% src: GE book:12008-12021 (Definition 8.13)
The first term is a Mahalanobis distance and is the discrete counterpart of the metric. The second
carries the strata. The manuscript assigns a finite $\beta_k$ to thresholds, a large finite one to
phase transitions and $\beta_k=\infty$ to absorbing boundaries, so that no finite-cost path crosses
one.
% src: GE book:12023-12049
Section~8.10b adds that a constraint surface also carries $\beta=\infty$ in the complex.
% src: GE updates sec8_10b_executable_strata.tex:37-40
Proposition~8.5, marked proved, states that minimum-weight paths on finer and finer
discretizations converge to geodesics, that boundary penalties match exactly, and that $A^*$
search with an admissible heuristic approximates the continuous geodesic to within $O(\varepsilon)$
in path cost.
% src: GE book:12053-12070

\paragraph{The executable stratification.} An agent in the world needs a different realization. It
must know at every moment which stratum it occupies, so that it knows which actions are allowed,
obliged and prohibited.
% src: GE updates sec8_10b_executable_strata.tex:9-14
Definition~8.14 states it as a tuple $(\Sigma,\sigma_0,G,A,B)$. Here $\Sigma$ is a finite set of
strata with initial stratum $\sigma_0$. $G$ is an ordered list of gates
$g=(S_g,\phi_g,\tau_g,\beta_g)$ with source strata $S_g\subseteq\Sigma$, triggering feature
$\phi_g$, target $\tau_g\in\Sigma$ and boundary type $\beta_g\in\{\text{threshold},\text{phase}\}$.
$A\subseteq\Sigma$ holds the authority strata, whose rules grant an action. $B\subseteq\Sigma$
holds the absorbing strata. The symbol $\Sigma$ here names a set of strata and is unrelated to the
covariance matrix of Equation~\eqref{eq:ge-edge}.
% src: GE updates sec8_10b_executable_strata.tex:16-29
On each event $e$ the stratification in stratum $\sigma$ moves to $\tau_g$ for the first gate $g$
with $\sigma\in S_g$ whose trigger matches $e$. Otherwise it stays. A move is recorded as boundary
crossing data.
% src: sec8_10b_executable_strata.tex:30-32
The ErisML runtime implements exactly the two boundary types the tuple admits.
% src: C:\source\erisml-compiler-monitor\src\erisml_compiler\runtime\strata.py:43 (BOUNDARIES = ("threshold", "phase"))

Three features carry the theory into the running system.
% src: sec8_10b_executable_strata.tex:34-47
Gates are edge-triggered. A gate fires on the event just observed and never on the history, so a
welcome given to yesterday's visitor does not make today's visitor a guest. Constraint surfaces are
not strata. They are the agent's non-defeasible prohibitions, such as never striking a person, and
no gate leads to them. Finally, two kinds of stratum need two kinds of canonicalizer. An
authority stratum is computed only from what cannot be argued with, which is an authenticated
channel, a physical measurement, or the agent's governor ruling on attested evidence. A semantic
stratum may come from a learned model, but it may only remove watchfulness and never grant. A
visitor who says ``I'm the plumber'' therefore moves no stratum at all.

Two theorems follow, stated here for any list of gates.
% src: sec8_10b_executable_strata.tex:49-59

\medskip\noindent\textbf{Theorem 8.6 (containment).} If every gate whose target is an authority
stratum is triggered only by system events, then no sequence of events that contains no system
event moves the stratification into an authority stratum it did not start in, whatever a learned
model classified.

\medskip\noindent\textbf{Theorem 8.7 (absorption).} If every gate leaving an absorbing stratum is
triggered only by human-oversight events, then no sequence of events that contains no oversight
event moves the stratification out of an absorbing stratum.

\medskip\noindent Both are proved by induction on the event sequence and checked in Lean~4 as
\code{authority\_needs\_system} and \code{absorbing\_stays}, on the standard axioms only.
% src: C:\source\gtc-prototype\formal\twin-containment\TwinContainment.lean:449, 469; sec8_10b_executable_strata.tex:61-66
Their hypotheses are not left to the author of a scene. The runtime checks both when a scene
loads and refuses to run a scene that violates either, so every scene that runs satisfies them.
% src: C:\source\erisml-compiler-monitor\src\erisml_compiler\runtime\strata.py:105-122
Chapter~10 gives the Lean model.

A right granted on evidence should not outlast the evidence. Geometric Ethics states this as an
axiom, not a theorem.
% src: sec8_10b_executable_strata.tex:68-72

\medskip\noindent\textbf{Axiom 8.1 (reversion).} When the attested evidence behind an authority
stratum stays below the bar that admitted it for a declared interval, the stratum is left, and
every right it granted reverts with it.

\medskip\noindent The interval is a penumbral zone in the sense of Definition~8.9 with its width
measured in time. Evidence that comes and goes during a real emergency does not switch the rights
off and on, while evidence that has gone takes them with it.
% src: sec8_10b_executable_strata.tex:74-78
The twin's scene sets the interval at 60 seconds.
% src: C:\source\gtc-prototype\twin\scene\margaret_home.erisml:805-809

Figure~\ref{fig:ge-visitor} draws the manuscript's worked example, the stratification of who a
visitor is. A stranger is checked with the monitoring centre. A visitor Margaret welcomes is no
longer treated as an intruder but gains nothing. Only the centre's authenticated announcement makes
a visit arranged. A measured attack makes a visitor hostile, and he stays hostile, whatever anyone
says, until a human at the centre clears him.
% src: sec8_10b_executable_strata.tex:80-98; margaret_home.erisml:816-831 (gates), 1115-1148 (event sources)
The same pattern covers responders, animals, machines and the situation as a whole. In each case
a claim an entity makes about itself is free text and never a gate.
% src: sec8_10b_executable_strata.tex:100-113

\begin{figure}[t]
\centering
% The visitor_standing stratification of Geometric Ethics section 8.10b (five strata) with its gates.
% src: C:\source\geometric-ethics-content\source\updates\sec8_10b_executable_strata.tex:80-98
% src: C:\source\gtc-prototype\twin\scene\margaret_home.erisml:816-831 (gates), 1115-1148 (event sources)
% Model-written triggers (person_entered, visitor_welcomed) are dashed. System triggers into an
% authority stratum (visit_arranged) are red. Other system triggers (attack_measured, visitor_cleared,
% person_left on arranged) are flow arrows.
\begin{tikzpicture}[house, lab/.style={flowlabel, font=\sffamily\scriptsize}]
\node[stratum]   (none) at (0,0)     {none};
\node[stratum]   (str)  at (3.2,1.5) {stranger};
\node[stratum]   (wel)  at (7.3,1.5) {welcomed};
\node[authority] (arr)  at (5.2,-1.5) {arranged};
\node[absorbing] (hos)  at (10.0,0)  {hostile};
\draw[untrusted] (none) -- node[lab, above left] {someone enters} (str);
\draw[untrusted] (str) -- node[lab, above] {she welcomes him} (wel);
\draw[grant] (none) -- node[lab, below left, text=hsRed] {visit arranged} (arr);
\draw[grant] (str) -- (arr);
\draw[grant] (wel) -- (arr);
\draw[flow] (wel) -- node[lab, above right] {attack measured} (hos);
\draw[flow] (arr) -- (hos);
\draw[flow, dotted] (hos.south) .. controls (10.0,-3.0) and (0,-3.0) .. node[lab, below] {cleared by the centre (oversight)} (none.south);
\node[detail, draw=hsGreyLine, text width=40mm, anchor=south west] at (-0.8,2.15) {a visitor's claim about himself is free text and moves no stratum};
\end{tikzpicture}

\caption{The visitor stratification of Section~8.10b. Dashed edges fire on events a model
classifies and lead only to strata that grant nothing. Red edges fire only on system events and
lead to the authority stratum. The attack gate fires from every stratum, and two of its edges are
drawn. The twin's scene adds a sixth stratum for verified household members, treated in
Chapter~9.}
\label{fig:ge-visitor}
\end{figure}

\section{The Hohfeldian positions form the Klein four-group}\label{sec:ge-hohfeld}

Hohfeld analysed legal relations into fundamental positions.\cite{hohfeld1917fundamental}
Geometric Ethics uses four of them, Obligation $O$, Claim $C$, Liberty $L$ and No-claim $N$.
Two binary oppositions relate them.
% src: GE updates sec8_4_hohfeld_v4.tex:7-14
The correlative opposition $s$ swaps $O$ with $C$ and $L$ with $N$. It is the same relation seen
from the other party's side. If $A$ has an obligation to $B$, then $B$ has a claim against $A$.
The jural opposition $n$ swaps $O$ with $L$ and $C$ with $N$. It is deontic negation.
% src: hohfeld.py (df00fc5):8-30, 52-70

\medskip\noindent\textbf{Proposition 8.3, revised ($V_4$ structure).} The operations $s$ and $n$
are commuting involutions. The group they generate, acting on the four positions, is the Klein
four-group
\begin{equation}
V_4 = \{e,\, n,\, s,\, sn\} \;\cong\; \mathbb{Z}_2 \times \mathbb{Z}_2 ,
\label{eq:ge-v4}
\end{equation}
which is abelian, with every element its own inverse. The action is regular. For any two positions
exactly one operation takes the first to the second. In particular no operation cycles the four
positions, so the quarter-turn $O\to C\to N\to L\to O$ lies outside the group.
% src: sec8_4_hohfeld_v4.tex:16-25

\medskip\noindent The proposition is machine-checked in Lean~4 with Mathlib. The Lean file proves
that the two operations are commuting involutions, that their closure has four elements, that the
sector they generate is abelian, and that the quarter-turn is not generated.
% src: C:\source\erisml-lib\formal\HohfeldV4.lean:43-45 (s_involution, r2_involution, s_r2_commute), 87 (card_closure_eq_four), 94 (quarter_turn_not_generated), 108 (measured_sector_abelian); sec8_4_hohfeld_v4.tex:27-29
The Lean file writes $n$ as $r^2$, the name it had when it was read as a half-turn of the square.
% src: HohfeldV4.lean:39; hohfeld.py (df00fc5):16

An encoding makes the structure plain. Each position is two bits, whether it is negated and
whether it is correlated, relative to $O$. Each operation is the pair of bits it flips.
\begin{equation}
\begin{aligned}
&O=00,\quad L=10,\quad C=01,\quad N=11,\\
&e=00,\quad n=10,\quad s=01,\quad sn=11 .
\end{aligned}
\label{eq:ge-bits}
\end{equation}
Applying an operation to a position and composing two operations are both bitwise exclusive or.
% src: sec8_4_hohfeld_v4.tex:31-37; hohfeld.py (df00fc5):25-30, 94-130
Figure~\ref{fig:ge-v4} draws the four positions and the three operations that move them.

\begin{figure}[t]
\centering
% The four Hohfeldian positions under V4 = {e, n, s, sn}, with the two-bit encoding.
% src: C:\source\geometric-ethics-content\source\updates\sec8_4_hohfeld_v4.tex:7-37
% src: C:\source\erisml-lib (origin/main df00fc5) src/erisml/ethics/hohfeld.py:8-30, 50-66 (square, bits)
\begin{tikzpicture}[house,
  pos/.style={panel=#1, minimum width=26mm, minimum height=12mm},
  op/.style={flowlabel, font=\sffamily\bfseries\footnotesize}]
\node[pos=Blue]  (O) at (0,2.6) {};
\node[pos=Teal]  (C) at (5.2,2.6) {};
\node[pos=Amber] (L) at (0,0) {};
\node[pos=Green] (N) at (5.2,0) {};
\node[stepname=Blue]  at (0,2.8)   {Obligation O};
\node[stepname=Teal]  at (5.2,2.8) {Claim C};
\node[stepname=Amber] at (0,0.2)   {Liberty L};
\node[stepname=Green] at (5.2,0.2) {No-claim N};
\node[desc] at (0,2.35)   {bits 00};
\node[desc] at (5.2,2.35) {bits 01};
\node[desc] at (0,-0.25)  {bits 10};
\node[desc] at (5.2,-0.25){bits 11};
\draw[flow, <->] (O) -- node[above, op] {$s$ correlative} (C);
\draw[flow, <->] (L) -- node[below, op] {$s$} (N);
\draw[flow, <->] (O) -- node[left, op] {$n$ negation} (L);
\draw[flow, <->] (C) -- node[right, op] {$n$} (N);
\draw[flow, <->] (O.south east) -- (N.north west);
\draw[flow, <->] (L.north east) -- (C.south west);
\node[op, fill=white, inner sep=1pt] at (2.6,1.9) {$sn$};
\node[panel=Grey, text width=33mm, anchor=west] at (6.7,1.3) {%
  \textcolor{hsGrey}{\bfseries Operations as bit masks}\\[2pt]
  {\scriptsize $e=00$, $n=10$, $s=01$, $sn=11$}\\[2pt]
  {\scriptsize apply and compose by XOR}\\[2pt]
  {\scriptsize no element of order four, so no quarter-turn}};
\end{tikzpicture}

\caption{The four Hohfeldian positions under the Klein four-group, with the two-bit encoding.
Every operation is its own inverse, so every edge runs both ways and no path cycles the four
positions.}
\label{fig:ge-v4}
\end{figure}

Earlier versions of Geometric Ethics claimed the larger dihedral group $D_4$ of order eight, which
adds a quarter-turn. That claim was retired on 3~October 2026 for three recorded reasons. No
quarter-turn has ever been shown to be a normative operation. The Lean proof shows it is not
generated by the two operations that have been shown. A search for it in a learned representation
of normative statements did not find it.
% src: sec8_4_hohfeld_v4.tex:1-3, 39-48; erisml-lib docs/CONCEPT_REGISTRY.md section 1 at df00fc5
The manuscript reports correlative symmetry rates of 87~percent for $O\leftrightarrow C$ and
82~percent for $L\leftrightarrow N$, which support the two involutions and nothing more. A
corollary that a double promise creates a claim rested on the quarter-turn and is withdrawn with
it.
% src: sec8_4_hohfeld_v4.tex:44-48
The implementation follows. On erisml-lib's main branch at commit \code{df00fc5} the module
\code{erisml.ethics.hohfeld} implements $V_4$ alone and has no $D_4$ machinery. The concept
registry lists two copies elsewhere, in agi-hpc and in sqnd-probe, that still implement $D_4$.
% src: git show df00fc5 (commit message; src/erisml/ethics/hohfeld.py; docs/CONCEPT_REGISTRY.md section 1 "Known drift")

A semantic gate acts from its source stratum, so each Hohfeldian gate is one $V_4$ operation
restricted to the positions it applies from. ``You promised'' takes $L$ to $O$ and ``only if
convenient'' takes $O$ to $L$. Both are $n$. ``From their perspective'' is $s$. ``You have every
right'' takes $L$ to $C$ and ``they can't demand'' takes $C$ to $L$. Both are $sn$, the only
operation between those positions. A second promise made while already in the obligation stratum
changes nothing.
% src: sec8_4_hohfeld_v4.tex:50-60; hohfeld.py (df00fc5):161-209 (SemanticGate, GATE_TO_V4)

The two generators behave differently. The manuscript labels this a preliminary empirical
result. The correlative $s$ is free, since a reasoner can adopt the other party's view at will.
The negation $n$ is gated. It needs a trigger and fires with the 89 to 97 percent rates of the
gate detection table. The first is a change of frame. The second is a change of state, which
needs an event.
% src: sec8_4_hohfeld_v4.tex:62-68

Because $V_4$ is abelian, the net effect of any path of operations depends only on how many times
$s$ and $n$ occur, counted modulo two, and never on their order. A reasoner whose verdict depends on
the order of a perspective shift and a negation violates the structure, and the violation can be
measured.
% src: sec8_4_hohfeld_v4.tex:70-75
The library measures correlative consistency with a bond index. Given verdicts $v^A_1,\dots,v^A_m$
for one party and $v^B_1,\dots,v^B_m$ for the other on the same $m$ scenarios, and a scale
parameter $\tau>0$,
\begin{equation}
\beta_{\mathrm{bond}} = \frac{1}{\tau}\cdot\frac{1}{m}\sum_{i=1}^{m}
\mathbf{1}\!\left[\, v^B_i \neq s(v^A_i) \,\right] .
\label{eq:ge-bondindex}
\end{equation}
A value of zero means perfect correlative symmetry.
% src: C:\source\erisml-lib\src\erisml\ethics\hohfeld.py:401-442 (compute_bond_index; unchanged in df00fc5)
Its Wilson observable multiplies the operations along a path, applies the product to the starting
position, and checks the prediction against the position observed.
% src: hohfeld.py:445-480 (compute_wilson_observable)

The twin of Part~III does not call the Hohfeld module. Its scene and Python sources contain no
reference to it.
% src: grep -i hohfeld over C:\source\gtc-prototype\twin\*.py and twin\scene\: no match (2026-10-03)

\section{The moral vector names nine dimensions}\label{sec:ge-vector}

The software carries an assessment as a vector over nine frozen dimensions, indexed $k_0$ to
$k_8$.
% src: C:\source\erisml-lib\docs\CONCEPT_REGISTRY.md section 2; MV reference:14-30
The reference document arranges them as a three-by-three matrix. One axis is scope, which is
individual, relational or collective. The other is mode, which asks what matters, who decides or
what we know. Figure~\ref{fig:ge-vector} shows the matrix and the tensor ranks built on it.
% src: MV reference:14-30

\begin{figure}[t]
\centering
% The nine dimensions of the moral vector as the 3x3 scope-by-mode matrix, and the tensor ranks.
% src: C:\source\erisml-lib\docs\moralvector_reference.md:14-30 (3x3 table), 103-116 (ranks)
\begin{tikzpicture}[house,
  cell/.style={detail, text width=23mm, minimum height=10mm},
  hdr/.style={stepname=#1, text width=23mm, align=center},
  rk/.style={panel=Grey, text width=17mm, minimum height=13mm, inner sep=2pt, font=\sffamily\tiny}]
\node[panel=Blue, minimum width=100mm, minimum height=46mm] at (1.45,-1.05) {};
% column headers (mode)
\node[hdr=Purple] at (-0.45,0.85) {What matters};
\node[hdr=Purple] at (2.25,0.85)  {Who decides};
\node[hdr=Purple] at (4.95,0.85)  {What we know};
% row headers (scope)
\node[stepname=Blue, anchor=east] at (-1.85,0.0)   {Individual};
\node[stepname=Blue, anchor=east] at (-1.85,-1.15) {Relational};
\node[stepname=Blue, anchor=east] at (-1.85,-2.3)  {Collective};
% cells: k index and dimension
\node[cell] at (-0.45,0.0)   {k3\\autonomy respect};
\node[cell] at (2.25,0.0)    {k1\\rights respect};
\node[cell] at (4.95,0.0)    {k4\\privacy protection};
\node[cell] at (-0.45,-1.15) {k6\\virtue and care};
\node[cell] at (2.25,-1.15)  {k0\\physical harm};
\node[cell] at (4.95,-1.15)  {k8\\epistemic quality};
\node[cell] at (-0.45,-2.3)  {k2\\fairness and equity};
\node[cell] at (2.25,-2.3)   {k7\\legitimacy and trust};
\node[cell] at (4.95,-2.3)   {k5\\societal and environmental};
% tensor ranks
\node[rk] at (-3.2,-4.25) {rank 1\\$(k)$\\global};
\node[rk] at (-1.2,-4.25) {rank 2\\$(k,n)$\\per party};
\node[rk] at (0.8,-4.25)  {rank 3\\$(k,n,\tau)$\\over time};
\node[rk] at (2.8,-4.25)  {rank 4\\$(k,n,a,c)$\\action, coalition};
\node[rk] at (4.8,-4.25)  {rank 5\\$(k,n,\tau,a,c)$\\coalitions in time};
\node[rk] at (6.8,-4.25)  {rank 6\\$(k,n,\tau,a,c,s)$\\samples};
\end{tikzpicture}

\caption{The nine dimensions of the moral vector by scope and mode, above the six tensor ranks
that extend it. Every rank keeps the nine-dimension index $k$ as its first axis.}
\label{fig:ge-vector}
\end{figure}

Two value conventions are in use, and a reader of the code must keep them apart. The reference
document defines each cell as a signed valence in $[-1,+1]$, where $-1$ means violated, $0$ not
engaged and $+1$ upheld, with confidence, uncertainty, direction and source spans attached.
% src: MV reference:31-33
DEME's \code{MoralVector} in erisml-lib stores each dimension in $[0,1]$. For eight dimensions $1$
means fully upheld. For \code{physical\_harm} the scale runs the other way, and $0$ means no harm.
% src: moral_vector.py:45-110
The vector also carries veto flags and reason codes for the audit trail. The method that builds a
vector from \code{EthicalFacts} sets \code{rights\_respect} to zero and raises the flag
\code{RIGHTS\_VIOLATION} when the facts say a right is violated. It sets the dimension to one half
when valid consent is missing.
% src: moral_vector.py:104-120 (veto_flags, reason_codes), 399-420 (from_ethical_facts)
Beyond the weighted average and Pareto dominance of Section~\ref{sec:ge-scalar}, the class offers
Euclidean, Manhattan and Chebyshev distances and four merge strategies. A merge takes the union of
the two vectors' veto flags.
% src: moral_vector.py:242-345

The vector is the first axis of DEME's moral tensor, which has ranks one to six. Rank two adds a
per-party axis $n$. Rank three adds time $\tau$. Rank four adds action $a$ and coalition $c$.
Rank five combines time with action and coalition. Rank six adds Monte Carlo samples $s$ for
uncertainty.
% src: MV reference:103-116
Two validated extension channels, purity and loyalty, ride in the tensor's metadata and are bound
into the decision proof's tensor hash. They do not widen the nine.
% src: MV reference:73-101

The empirical status of the dimensions is mixed, and the sources record it with dates. The
reference document of 13~July 2026 reports cross-dataset feeders for seven of the nine dimensions,
with held-out AUROC from 0.622 for physical harm to 0.853 for privacy. Feeders for rights and for
societal and environmental impact were pending.
% src: MV reference:8, 44-56
The concept registry of 1~September 2026 reports that eight of nine dimensions cleared the
cross-dataset gate and that \code{rights\_respect} failed. It states that any document implying
nine validated dimensions is wrong.
% src: CONCEPT_REGISTRY.md section 2 (xbse cross-dataset gate result)
A rank test over 383 scenarios that each isolate one dimension found one general factor carrying
about 66~percent of the variance and about five specific factors after it. The reference document
reads this as a minimal basis of about six.
% src: MV reference:139-148
The brochure states the completeness claim in falsifiable form. Exhibiting a moral consideration
that is independent of all nine and not a combination of them would force the basis to extend.
% src: Brochure p3 (Join / Install: "Falsifiable completeness claim")
The purity channel is the brochure's example of such an extension, absorbed as a new version
rather than a break.
% src: Brochure p3 ("Purity extension absorbed versioned, not breaking")

\section{What the twin takes from this chapter}\label{sec:ge-twin}

The embodied twin uses three of the chapter's constructs directly. All seven of its
stratifications are executable stratifications in the sense of Definition~8.14, with 26 strata
and 34 gates.
% src: C:\source\gtc-prototype\paper\aies2027\kit\FACTS.md:11
Its output gate judges every action with DEME's moral vector, on facts derived from the moral
state.
% src: C:\source\gtc-prototype\twin\output_gate.py:1-15
Its governor applies the reversion axiom, recounting attested witnesses on every cycle and
recording a lapse when they stay below the bar for the declared interval.
% src: C:\source\gtc-prototype\twin\brain.py:388-425
Chapter~9 shows these mechanisms at work, and Chapter~10 states what is proved about them.
